\par
\noindent \begin{minipage}[t]{\columnwidth}
\vspace{0pt}
\section{Computational cost}
\label{sec:cost}
To test feasibility with lower compute, each planning run uses a single RTX 3090. Native and ECP share a fixed checkpoint and tasks; paired ten-task batches run sequentially on the same GPU with alternating order. The implementation uses PyTorch 2.3.0 and CUDA 12.1.\par\smallskip
\normalsize
GPU-synchronized timers include CUDA Graph capture and simulator replay, excluding model loading and the separate physics audit. \Cref{tab:cost} shows that fewer replans offset the per-replan overhead and reduce mean task time. Peak allocated memory stays below one GiB.\par

\end{minipage}\hfill
\begin{minipage}[t]{\columnwidth}
\vspace{0pt}
\par\smallskip
{\centering
\centering
\captionof{table}{\textbf{Planning time and memory on PushObj.} 350 paired tasks, fixed ResNet. Task time includes failures capped at 30 replans. P95: 95th percentile; memory: maximum PyTorch peak.}
\label{tab:cost}
\small
\setlength{\tabcolsep}{3.5pt}
\begin{tabular*}{\columnwidth}{@{\extracolsep{\fill}}lrr@{}}
\toprule
Metric & Without ECP & With ECP\\
\midrule
Success (\%) &61.14&72.29\\
Mean replans / task &15.59&13.66\\
Planning (ms / replan) &747.24&768.27\\
GD (ms / replan) &747.14&752.39\\
Mean task time (s) &15.25&13.41\\
P95 task time (s) &31.60&32.49\\
Total planning time (s) &4078.45&3672.32\\
Peak allocated memory (GiB) &0.823&0.967\\
Peak reserved memory (GiB) &1.342&1.469\\
\bottomrule
\end{tabular*}

\par}

\end{minipage}\par
